\chapter{Contexte et préliminaires}
\label{ch:background}

% Introduisez uniquement les notations et résultats antérieurs utilisés
% ensuite. Indiquez clairement ce qui est classique et citez sa source ; ne
% transformez pas ce chapitre en panorama sans lien avec la suite.

\begin{definition}[Graphe orienté fini]
\label{def:directed-graph}
Un \emph{graphe orienté fini} est un couple $G=(V,E)$, où $V$ est un ensemble
fini de sommets et $E\subseteq V\times V$ un ensemble d'arêtes orientées.
\end{definition}

\begin{definition}[Accessibilité]
\label{def:reachability}
Soit $G=(V,E)$ un graphe orienté fini. Un sommet $v\in V$ est
\emph{accessible} depuis $u\in V$, ce que l'on note $u\mathrel{R^*}v$, s'il
existe un chemin de longueur $k\geq 0$ de $u$ vers $v$. Un chemin de longueur
nulle rend chaque sommet accessible depuis lui-même.
\end{definition}

En notant $R=E$ la relation d'arêtes et $I_V$ la relation identité sur $V$, la
clôture réflexive transitive est
\begin{equation}
  R^* = \bigcup_{k\geq 0} R^k,
  \label{eq:transitive-closure}
\end{equation}
où $R^0=I_V$ et les puissances de relations sont définies par composition. Cette
construction classique est présentée sous des angles complémentaires dans
\cite{diestel2017,cormen2022}.

\begin{theorem}[Transitivité de l'accessibilité]
\label{thm:transitivity}
Pour tous sommets $u,v,w\in V$, si $u\mathrel{R^*}v$ et
$v\mathrel{R^*}w$, alors $u\mathrel{R^*}w$.
\end{theorem}

\begin{proof}
La première hypothèse fournit un chemin de $u$ à $v$, et la seconde un chemin
de $v$ à $w$. Leur concaténation forme un chemin de $u$ à $w$. Ainsi,
$u\mathrel{R^*}w$ d'après la \cref{def:reachability}.
\end{proof}

Le \cref{thm:transitivity} découle aussi directement de
l'\cref{eq:transitive-closure} : l'union contient les chemins de toute longueur
finie et reste fermée par concaténation.

\begin{corollary}[L'accessibilité est un préordre]
\label{cor:preorder}
La relation $R^*$ est réflexive et transitive sur $V$ : la réflexivité vient des
chemins de longueur nulle (\cref{def:reachability}) et la transitivité est le
\cref{thm:transitivity}.
\end{corollary}
